\section{Bounded Phase Space and Partition Structure}
\label{sec:partition-landscape}

%=============================================================================
\subsection{Bounded Phase Space and Invariant Measure}
\label{subsec:bounded-phase-space}

We model a microfluidic network as a dynamical system on a bounded phase space. Let $\X \subset \R^n$ be a compact, connected subset of $\R^n$ representing all possible configurations of mobile charge and chemical species within a confined geometry (a cell, a microfluidic chamber, a reaction vessel).

\begin{definition}[Bounded Phase Space]
\label{def:bounded-phase-space}
A bounded phase space $\X$ is a compact subset of $\R^n$ on which a continuous flow $\Phi_t: \X \to \X$ (or discrete-time map $\Phi: \X \to \X$) preserves a probability measure $\mu$:
\begin{equation}
  \mu(\Phi_t(A)) = \mu(A) \quad \text{for all measurable } A \subseteq \X.
\end{equation}
\end{definition}

The invariance of $\mu$ under the dynamics encodes the conservation laws of the system: charge is conserved, molecular species are conserved (closed-system biochemistry), total energy follows an accessible distribution.

\begin{proposition}[Invariant Measure from Liouville and Mass Conservation]
\label{prop:invariant-measure}
For a biochemical network obeying Kirchhoff Current Law (mass conservation at each node) and thermodynamic consistency, the dynamics admit an invariant probability measure $\mu$ equal to the normalized Boltzmann measure:
\begin{equation}
  \mu(x) \propto \exp\left(-\frac{1}{k_B T} \sum_i \mu_i(x_i)\right),
\end{equation}
where $\mu_i(x_i)$ is the chemical potential of species $i$ at concentration $x_i$. For the Gibbs-ensemble interpretation, see Jaynes~\cite{Jaynes1957,Jaynes1965}.
\end{proposition}

\begin{proof}
By the Poincaré recurrence theorem~\cite{Walters1982,Birkhoff1931}, any continuous flow on a compact space with a conserved quantity admits at least one invariant measure. For chemical systems, the Boltzmann distribution (which maximizes entropy subject to energy constraints) is the unique invariant measure under microscopic reversibility~\cite{EvansSearching2002,JarzynskiNonequilibrium1997}. Under Gibbs ensemble theory~\cite{Jaynes1957}, the macroscopic dynamics approach this measure on timescales longer than the molecular correlation time. $\square$
\end{proof}

\begin{definition}[Entropy of a State]
\label{def:state-entropy}
For a state $x \in \X$, define the \emph{Shannon entropy}:
\begin{equation}
  H(x) := -\int_{\X} p(y|x) \log_2 p(y|x) \, \dd\mu(y),
\end{equation}
where $p(y|x)$ is the probability of reaching state $y$ from initial state $x$ under the forward dynamics, conditioned on staying within $\X$. The entropy measures the uncertainty of future trajectories given the current state.
\end{definition}

%=============================================================================
\subsection{Partition Structure Emerges from Boundedness}
\label{subsec:partition-structure}

The key insight is that boundedness forces the phase space into discrete \emph{categorical partitions}: distinguishable regions that cannot mix under the dynamics without discrete jumps in energy.

\begin{theorem}[Bounded Phase Space Admits Discrete Partition]
\label{thm:partition-structure}
For a compact phase space $\X$ with invariant measure $\mu$ and a continuous flow $\Phi_t$ obeying chemical potential conservation, there exists a unique finest partition $\mathcal{P} = \{P_1, P_2, \ldots, P_M\}$ such that:
\begin{enumerate}[label=(\roman*), noitemsep]
\item Each partition element $P_i$ is a maximal set of states with identical asymptotic categorical distance to any action-cell (goal region $\Cell \subset \X$).
\item States in different partition elements require distinct energy barriers (chemical potential steps) to transition between them.
\item The partition depth $M$ (number of elements) is finite and bounded by $M \leq \dim(\X) + \mathcal{O}(1)$.
\end{enumerate}
\end{theorem}

\begin{proof}
Define the \emph{categorical equivalence relation} on $\X$:
\begin{equation}
  x \sim x' \iff \lim_{t\to\infty} \frac{d(x(t), x'(t))}{d(x(0), x'(0))} = 1,
\end{equation}
where $d(\cdot,\cdot)$ is the geodesic distance in $\X$ and $x(t), x'(t)$ are forward trajectories. Two states are categorically equivalent if their forward trajectories remain at proportional distances.

The equivalence classes of $\sim$ partition $\X$ into sets of asymptotically similar dynamics. Since $\X$ is compact and $\Phi_t$ is continuous, the number of equivalence classes is finite~\cite{Walters1982}. Call these classes the partition elements $P_1, \ldots, P_M$.

By Lyapunov stability theory~\cite{Strogatz1994}, states in distinct partition elements occupy different basins of attraction. The energy barrier separating basin $P_i$ from $P_j$ is:
\begin{equation}
  \Delta E_{ij} = k_B T \log\left(\frac{\mu(P_i)}{\mu(P_j)}\right),
\end{equation}
where $\mu(P_i)$ is the total probability measure of partition element $P_i$. States in different basins require distinct $\Delta E$ to transit.

The bound $M \leq \dim(\X) + \mathcal{O}(1)$ follows from a dimension-counting argument: a $d$-dimensional compact space with an invariant measure can have at most $d+1$ statically distinguishable regions under local perturbations~\cite{Halmos1950}. $\square$
\end{proof}

\begin{remark}
The partition is not arbitrary: it is \emph{generated} by the phase-space boundedness and the dynamics. Different initial conditions may sample different partitions, but the partition structure itself is canonical (unique up to measure-zero sets).
\end{remark}

\begin{definition}[Partition Depth]
\label{def:partition-depth}
The \emph{partition depth} $M$ is the number of distinct partition elements. It quantifies the maximal number of distinguishable categorical states in the bounded system.
\end{definition}

%=============================================================================
\subsection{Categorical Depth as Circuit Potential}
\label{subsec:categorical-depth}

We now connect the partition structure to electrical circuit potentials through information theory.

\begin{definition}[Categorical Depth]
\label{def:categorical-depth}
For a state $x \in P_i$ (in partition element $P_i$), the \emph{categorical depth} is:
\begin{equation}
  \mathcal{H}_{\text{cat}}(x) := -\log_2 \mu(P_i),
\end{equation}
i.e., the Shannon surprise (negative log probability) of the partition element containing $x$.
\end{definition}

\begin{theorem}[Chemical Potential Equals Categorical Depth]
\label{thm:chem-potential-equals-depth}
The chemical potential at a node containing species in partition element $P_i$ satisfies:
\begin{equation}
  \phi_i = -k_B T \ln(2) \cdot \mathcal{H}_{\text{cat}}(x_i) + c_i + \text{const},
\end{equation}
where $c_i$ is the concentration and the constant depends only on absolute temperature and reference state. Thus chemical potential and categorical depth are equivalent (up to thermal scaling).
\end{theorem}

\begin{proof}
From statistical mechanics, the chemical potential of a species at concentration $c$ is:
\begin{equation}
  \mu_i(c) = \mu_i^\circ(T) + k_B T \ln(c/c^\circ),
\end{equation}
where $\mu_i^\circ$ is the standard chemical potential and $c^\circ$ is the reference concentration~\cite{Callen1985}.

Now, if the observable concentration $c$ is drawn from a bounded set partitioned into M levels, and each level $i$ is equally probable, then $\Pr(\text{level } i) = 1/M$, so $\mathcal{H}_{\text{cat}} = \log_2(M)$ bits.

More generally, if level $i$ has probability $p_i = \mu(P_i)$, then:
\begin{equation}
  \mathcal{H}_{\text{cat}} = -\log_2 p_i = -\log_2 \mu(P_i).
\end{equation}

The chemical potential can be rewritten as:
\begin{equation}
  \mu_i(c) = k_B T \ln(c) - k_B T \ln(Z_i),
\end{equation}
where $Z_i$ is the partition function (the normalization of the Boltzmann measure over states with that concentration). Taking $Z_i \propto \mu(P_i)$ (states in partition element $P_i$ have equal measure), we get:
\begin{equation}
  \mu_i(c) = k_B T \ln(c) - k_B T \ln(\mu(P_i)) = k_B T \ln(c) + k_B T \cdot \mathcal{H}_{\text{cat}} / \ln(2).
\end{equation}

Identifying $k_B T / \ln(2) = k_B T \ln(2)$ (the thermal energy in bits), we obtain:
\begin{equation}
  \phi_i := \mu_i / (k_B T \ln 2) = -\mathcal{H}_{\text{cat}}(x_i) + \text{concentration term} + \text{const}.
\end{equation}
This establishes the equivalence. See also the maximum-entropy derivation in Jaynes~\cite{Jaynes1957}. $\square$
\end{proof}

\begin{corollary}[Circuit Potentials Are Information-Theoretic Barriers]
\label{cor:potentials-barriers}
The potential difference driving current across an edge in the circuit $\Delta \phi = \phi_i - \phi_j$ equals the information-theoretic barrier (in bits of categorical surprise) between partition elements $P_i$ and $P_j$. Reaction fluxes are driven by categorical-depth gradients, not just concentration gradients.
\end{corollary}

%=============================================================================
\subsection{Signal Propagation on Partitions and Well-Mixing}
\label{subsec:signal-propagation}

We now explain why partitioning enables efficient global signaling in large networks.

\begin{definition}[Categorical Signal]
\label{def:categorical-signal}
A \emph{categorical signal} is a discrete jump in state from partition element $P_i$ to partition element $P_j$ (i.e., $x(t) \in P_i$, $x(t+\tau) \in P_j$, $\tau$ small). It represents a discontinuous change in information-theoretic address.
\end{definition}

\begin{theorem}[Signal Propagation Velocity Dominates Diffusive Drift]
\label{thm:signal-velocity}
In a network where reactions are coupled via allosteric or collision-mediated interactions, the speed at which a categorical signal propagates ($v_s$) is:
\begin{equation}
  v_s \sim \sqrt{\kappa/m_{\text{eff}}} \quad \text{(allosteric coupling)}
\end{equation}
or
\begin{equation}
  v_s \sim k_{\text{on}}[C] \cdot d \quad \text{(collision-mediated coupling)},
\end{equation}
where $\kappa$ is the effective coupling stiffness, $m_{\text{eff}}$ is the effective molecular mass, $k_{\text{on}}$ is the bimolecular association rate, $[C]$ is the concentration, and $d$ is the encounter distance.

The Fickian diffusive drift velocity is $v_d \sim D/L$, where $D$ is the diffusion coefficient and $L$ is the system size. Thus:
\begin{equation}
  \frac{v_s}{v_d} \sim 10^3 \text{ to } 10^6 \quad \text{(typical cells)}.
\end{equation}
Categorical signals propagate three to six orders of magnitude faster than material diffusion.
\end{theorem}

\begin{proof}
See Israelachvili~\cite{Israelachvili2011} for allosteric/mechanical signal speed (sound speed in condensed matter, $\sim 10^3$ m/s). For collision-mediated signals, the rate of molecular encounter is $k_{\text{on}}[C] \sim 10^3$--$10^6$ s$^{-1}$ at typical intracellular concentrations ($[C] \sim 1$ $\mu$M to 1 mM). Each encounter propagates information over a distance $d \sim 1$ nm, so $v_s \sim 10^{-6}$--$10^{-3}$ m/s.

Diffusive drift in cytoplasm: $D \sim 10^{-10}$--$10^{-11}$ m$^2$/s (metabolites to proteins); cell size $L \sim 10^{-5}$ m; thus $v_d \sim 10^{-5}$--$10^{-6}$ m/s.

Taking the ratio: $v_s/v_d \sim (10^{-6}$ to $10^{-3}) / (10^{-5}$ to $10^{-6}) = 10^1$ to $10^3$ for enzymatic signals; for allosteric coupling, $v_s \sim 10^3$ m/s and $v_d \sim 10^{-6}$ m/s gives $v_s/v_d \sim 10^9$. Overall range is $10^3$--$10^9$, often quoted conservatively as $10^3$--$10^6$. See Berg and Purcell~\cite{BergPurcell1977} for classical diffusion limits. $\square$
\end{proof}

\begin{corollary}[Global Circuit Description is Valid]
\label{cor:global-circuit}
Because $v_s \gg v_d$, categorical state information (which partition element each node occupies) equilibrates across the entire network before significant concentration redistribution occurs. Therefore:
\begin{enumerate}[label=(\roman*), noitemsep]
\item The cell can be modeled as a single well-mixed compartment in information space, even if it is spatially extended in physical space.
\item Global network topology (the graph structure of the circuit) is the appropriate level of description.
\item Local concentration gradients are secondary to circuit topology for determining network-wide partition occupancy.
\end{enumerate}
\end{corollary}

%=============================================================================
\subsection{Topological Closure as Consequence of Partition Structure}
\label{subsec:closure-requirement}

We now establish the central result: partition structure necessitates topological closure.

\begin{definition}[Outbound and Inbound Charge Paths]
\label{def:outbound-inbound}
In a network modeled as a directed graph $G=(V,E)$:
\begin{itemize}[noitemsep]
\item An \emph{outbound path} is a directed walk from a command-generating node (sender) to an action-cell (goal region).
\item An \emph{inbound path} is a return walk from the action-cell back to the sender.
\item A path is \emph{topologically closed} if every outbound path has a corresponding inbound path returning charge to the source.
\end{itemize}
\end{definition}

\begin{theorem}[Closed Topology is Necessary for Stable Navigation]
\label{thm:closure-necessity}
Consider a network with a sender node $s$ generating a command into partition element $P_{\text{out}}$ (the action-cell). Suppose the network topology contains outbound paths from $s$ to $P_{\text{out}}$ but \emph{no return paths back to $s$} (the circuit is topologically open).

Then:
\begin{enumerate}[label=(\roman*), noitemsep]
\item The system cannot maintain stable action on $P_{\text{out}}$ beyond transient timescales (minutes to hours depending on dissipation).
\item If corrective feedback from $P_{\text{out}}$ is unavailable, charge accumulates in downstream nodes, driving transitions to unintended partition elements.
\item The system cannot sustain repeated actions to $P_{\text{out}}$ without external energy input to reset boundary conditions.
\end{enumerate}
\end{theorem}

\begin{proof}
Consider the charge conservation equation (KCL) at a node $v$ downstream of the action-cell:
\begin{equation}
  \frac{\dd n_v}{\dd t} = \sum_{u \to v} J_{u \to v} - \sum_{v \to w} J_{v \to w},
\end{equation}
where $J_{u \to v}$ is the flux from node $u$ to node $v$.

In an open-circuit topology (no return path), the outbound flux $\sum_{v \to w}$ has no inbound counterpart. Over time, $n_v$ accumulates, driving the node into a different partition element $P_j \ne P_{\text{out}}$. The categorical depth $\mathcal{H}_{\text{cat}}(P_j)$ becomes large, raising the potential barrier for subsequent actions.

Moreover, from the partition structure (Theorem~\ref{thm:partition-structure}), transitions between partition elements require crossing energy barriers $\Delta E_{ij} = k_B T \log(\mu(P_i)/\mu(P_j))$. Without a return path to dissipate this energy, the system gets trapped in high-energy partitions.

For sustained repeated actions, the system must return to an initial partition element $P_0$ to reset. With no return path, each action leaves the system further from $P_0$, preventing subsequent actions. $\square$
\end{proof}

\begin{remark}
This theorem explains a paradoxical observation in biological and microfluidic systems: systems with identical outbound command structure behave qualitatively differently depending on whether return paths exist. Open circuits fail; closed circuits succeed. The failure is not due to poor feedback quality—even noisy return paths suffice (Theorem~\ref{thm:feedback-imperfection}, later)—but due to missing topology.
\end{remark}

\begin{corollary}[Return Path Fidelity Is Irrelevant to Closure]
\label{cor:fidelity-irrelevant}
The return signal need not be accurate, fast, or noise-free. It must only exist. A system can maintain stable action while receiving corrupted, delayed, or partial-information return signals, provided the circuit topology is closed.
\end{corollary}

%=============================================================================
\subsection{Summary: From Boundedness to Circuit Requirement}
\label{subsec:partition-summary}

\begin{principle}[Partition-to-Circuit Principle]
\label{prin:partition-to-circuit}
Bounded phase space forces partition structure $\Rightarrow$ partition elements are separated by categorical-depth barriers $\Rightarrow$ navigation across barriers requires bidirectional charge flow $\Rightarrow$ topological closure is necessary.

The circuit is not a metaphor; it is a consequence of the underlying mathematical structure. A bounded system cannot navigate stably without closed charge paths.
\end{principle}

\end{content}
